\chapter{Projektverlauf}\label{ch:verlauf}

Das Projekt lief von der ersten Zeile Code am 10.\,Juni 2026 bis zur Abschlussauswertung am 2.\,Oktober 2026, also über 17 Kalenderwochen. Rückblickend lässt es sich in vier Phasen gliedern, die sich jeweils durch ein dominierendes Problem auszeichnen (\cref{fig:zeitleiste}).

\begin{figure}[htbp]\centering
  \resizebox{\linewidth}{!}{\zeitleiste}
  \caption{Zeitleiste des Projekts mit Phasen und Meilensteinen (Kalenderwochen 24--40, 2026). Die Phasen sind in den folgenden Abschnitten beschrieben.}\label{fig:zeitleiste}
\end{figure}

\section{Phasen im Überblick}

\paragraph{Phase 1: Prototyp und AWS-Grundlage (KW\,24--27, Juni bis Anfang Juli).}
Ein Scrapy-Spider, der PDF- und Word-Dokumente von einer Start-URL herunterlädt, wird um einen \texttt{JobRunner} (ein Job je URL, parallel begrenzt) und einen persistenten Visited-Store erweitert. Das erste Architekturdiagramm entsteht (Architektur v1, \cref{sec:arch-v1}). Danach folgt die Cloud-Anbindung: Lambda-Container-Image in ECR, S3 für Dokumente, DynamoDB für den Visited-Store, anschließend Webapp mit Authentifizierung auf einer EC2-Instanz. Dominierendes Thema: \emph{überhaupt etwas Lauffähiges in die Cloud bringen}.

\paragraph{Phase 2: Klassifikator und Plattform (KW\,28--30, Juli).}
Der Endlosschleifen-Bug bei Mehrfach-Formaten wird behoben, Jobs laufen asynchron mit Fortschrittsanzeige, ein Terraform-Stack soll die Infrastruktur beschreiben. Der erste TF-IDF-Klassifikator wird trainiert und in die Scraping-Pipeline integriert. Es folgen die Human-in-the-Loop-Feedbackschleife, CI/CD über GitHub Actions und, als Antwort auf das 15-Minuten-Limit von Lambda, der Fargate-Bulk-Pfad (Architektur v2, \cref{sec:arch-v2}). Dominierendes Thema: \emph{Qualität und Skalierung des Crawls}.

\paragraph{Phase 3: Großlauf und Recall-Analyse (KW\,32--36, August bis Anfang September).}
Am 12.\,August läuft die AWS-Sandbox ab, alle Daten sind verloren; die Umgebung wird über ein Bootstrap-Skript neu aufgebaut, und noch am selben Tag wird der Großlauf über 411 Hochschulen in der neuen Umgebung wiederholt. Die Auswertung zeigt, dass nur ein Teil der erwarteten Handbücher gefunden wurde. Es entstehen das Evaluationsmodul, die Discovery (Common Crawl, Serper, Firecrawl), der OCR-Fallback, die LLM-Zweitstufe und der Seeds-only-Modus (Architektur v3, \cref{sec:arch-v3}). Dominierendes Thema: \emph{Recall}.

\paragraph{Phase 4: Deep Run und LLM-Bereinigung (KW\,38--40, September bis Oktober).}
Nach der Budget-Sperre der Sandbox wird der Crawl um Breitensuche und Laufzeit-Wachhund erweitert und als zwölfstündiger Deep Run gestartet (\zDocsTotal{} PDFs kumuliert). Weil die TF-IDF-Positive in großer Zahl falsch sind, werden rund \zLLMReviewed{} Dokumente von einem LLM geprüft, Scans per OCR nachbearbeitet und die Abschlussauswertung mit 23 Grafiken erstellt. Dominierendes Thema: \emph{Precision}.

\section{Wochenweise Aktivität}
\Cref{tab:wochen-jan} schlüsselt die Aktivität von Jan Soetebeer wochenweise auf. Die Aufgaben stammen aus der Commit-Historie (28 Commits auf \texttt{main}, 42 inklusive Nebenzweigen) und den Entwicklungssitzungen. Die Stundenangaben sind \emph{Schätzungen}; sie wurden anhand der dokumentierten Sitzungs- und Commit-Zeiten sowie des Aufwands für Warte- und Betriebszeiten (Läufe, AWS-Konsole, Recherche) angesetzt \TODO{Stunden von Jan prüfen und anpassen (Quelle: bericht/daten/stunden.csv)}.

{\small
\begin{longtable}{@{}l l p{8.9cm} r@{}}
\caption{Wochenübersicht Jan Soetebeer (Stunden geschätzt).}\label{tab:wochen-jan}\\
\toprule
\textbf{Woche} & \textbf{Zeitraum} & \textbf{Tätigkeiten} & \textbf{Std.} \\
\midrule
\endfirsthead
\toprule
\textbf{Woche} & \textbf{Zeitraum} & \textbf{Tätigkeiten} & \textbf{Std.} \\
\midrule
\endhead
\bottomrule
\endlastfoot
% AUTOMATISCH ERZEUGT aus daten/stunden.csv
KW\,24 & 08.06.--14.06. & Projektstart; Scrapy-Prototyp (DocumentSpider, PDF/Word-Download); JobRunner mit parallelen Jobs; Visited-Store; erstes Architekturdiagramm (Architektur v1) & ca.\,10 \\
KW\,25 & 15.06.--21.06. & Einarbeitung AWS (Lambda, ECR, IAM); keine Commits [Jan: prüfen] & ca.\,2 \\
KW\,26 & 22.06.--28.06. & AWS-Integration: Lambda-Container-Image, ECR, S3-Pipeline, DynamoDB-Visited-Store, IAM-Policies; Frontend-Grundgerüst (Login, Tabs) & ca.\,12 \\
KW\,27 & 29.06.--05.07. & Authentifizierung (bcrypt, Session-Token); Hosting der Webapp auf EC2 mit Docker Compose und Caddy; Rollen- und Modellverwaltung & ca.\,10 \\
KW\,28 & 06.07.--12.07. & Endlosschleifen-Bug (Timeout bei Mehrfach-Formaten) behoben; asynchrone Jobs mit Fortschrittsanzeige; Terraform-Stack; Tiefen-Crawl; Ergebnis- und Review-UI; erster Klassifikator (TF-IDF + logistische Regression, Gruppen-Kreuzvalidierung); Scraper-Integration & ca.\,16 \\
KW\,29 & 13.07.--19.07. & MVP-Präsentation erstellt (15.07.-16.07.); Deployment-Versuche Lambda-Image und EC2; Klassifikator in Lambda-Image eingebettet & ca.\,8 \\
KW\,30 & 20.07.--26.07. & Human-in-the-Loop-Feedbackschleife (Review-UI, S3-Feedback, Nachtraining); CI/CD mit GitHub Actions (OIDC, SSM); CSV-Import der Hochschulliste; Skalierungs-Roadmap; fokussierter Crawl (Best-First, Sitemap, Subdomains); Extraktionsprofile; Fargate-Bulk-Pfad (Architektur v2) & ca.\,22 \\
KW\,31 & 27.07.--02.08. & keine dokumentierte Aktivität & ca.\,0 \\
KW\,32 & 03.08.--09.08. & Planung der Evaluation des ersten Großlaufs & ca.\,3 \\
KW\,33 & 10.08.--16.08. & Ablauf der Sandbox: kompletter Neuaufbau der AWS-Umgebung (Bootstrap-Skript); erster Großlauf über 411 Hochschulen (6,3 h); Evaluationsmodul mit HTML-Bericht & ca.\,14 \\
KW\,34 & 17.08.--23.08. & Recall-Analyse (145 Hochschulen ohne Treffer); Schwellenwert-Anpassung; Discovery-Modul (Common Crawl); OCR-Fallback & ca.\,8 \\
KW\,35 & 24.08.--30.08. & Serper- und Firecrawl-Anbindung; LLM-Zweitstufe auf Amazon Bedrock; End-to-End-Test auf Beispielhochschulen & ca.\,9 \\
KW\,36 & 31.08.--06.09. & Webapp-Schalter für Discovery/Rendering; parallele Bedrock-Aufrufe; Nachtraining mit 5.000+ Modulhandbüchern; Seeds-only-Modus; Per-Hochschule-Diagnose; Seeds-only-Lauf (+43 Hochschulen) & ca.\,12 \\
KW\,37 & 07.09.--13.09. & keine dokumentierte Aktivität (abwesend) & ca.\,0 \\
KW\,38 & 14.09.--20.09. & Budget-Sperre der Sandbox analysiert; Robots-Bypass für Discovery-Seeds; Fehler in der Evaluation (LLM-Urteile überschrieben) behoben; Breadth-Fairness und Watchdog; Deep Run (12 h, 132.193 PDFs) & ca.\,14 \\
KW\,39 & 21.09.--27.09. & keine dokumentierte Aktivität & ca.\,0 \\
KW\,40 & 28.09.--04.10. & LLM-Review von rund 27.700 Dokumenten; OCR von 3.553 Scans; Pilotstichprobe; Abschlussauswertung mit 23 Grafiken; Projektbericht & ca.\,16 \\
\midrule
\multicolumn{3}{l}{\textbf{Summe (Schätzung)}} & \textbf{ca.\,156} \\

\end{longtable}}

\Cref{fig:B01_stunden_pro_woche} zeigt dieselben Werte als Säulendiagramm. Die Spitzen liegen in KW\,30 (CI/CD, Feedback-Loop, fokussierter Crawl und Fargate an zwei Tagen), KW\,28 (Klassifikator und Review-UI) sowie in KW\,33 (Sandbox-Neuaufbau und erster Großlauf). Die Pausen in KW\,31, 37 und 39 sind Urlaubs- bzw. Abwesenheitszeiten, in denen nichts gelaufen ist; das ist für die Budget-Sperre in Kapitel~\ref{ch:aws} relevant. Insgesamt ergeben sich rund \zStundenJan{}~Stunden (Phase~1: \zStundenPhaseEins{}\,h, Phase~2: \zStundenPhaseZwei{}\,h, Phase~3: \zStundenPhaseDrei{}\,h, Phase~4: \zStundenPhaseVier{}\,h).

\abb[0.92]{htbp}{B01_stunden_pro_woche}{Geschätzter Arbeitsaufwand von Jan Soetebeer je Kalenderwoche, eingefärbt nach Projektphase.}

\subsection*{Wochenübersicht Sascha Lauk}
\saschaPH{Wochenweise Aktivität und Stunden ergänzen. Die Tabelle ist vorbereitet; Zeitraum und Wochen sind identisch zu Tabelle~\ref{tab:wochen-jan}.}

{\small
\begin{longtable}{@{}l l p{8.9cm} r@{}}
\caption{Wochenübersicht Sascha Lauk (Platzhalter).}\label{tab:wochen-sascha}\\
\toprule
\textbf{Woche} & \textbf{Zeitraum} & \textbf{Tätigkeiten} & \textbf{Std.} \\
\midrule
\endfirsthead
\toprule
\textbf{Woche} & \textbf{Zeitraum} & \textbf{Tätigkeiten} & \textbf{Std.} \\
\midrule
\endhead
\bottomrule
\endlastfoot
% AUTOMATISCH ERZEUGT aus daten/stunden.csv - Platzhalter für Sascha Lauk
KW\,24 & 08.06.--14.06. & \saschaCell{} & \saschaCell{} \\
KW\,25 & 15.06.--21.06. & \saschaCell{} & \saschaCell{} \\
KW\,26 & 22.06.--28.06. & \saschaCell{} & \saschaCell{} \\
KW\,27 & 29.06.--05.07. & \saschaCell{} & \saschaCell{} \\
KW\,28 & 06.07.--12.07. & \saschaCell{} & \saschaCell{} \\
KW\,29 & 13.07.--19.07. & \saschaCell{} & \saschaCell{} \\
KW\,30 & 20.07.--26.07. & \saschaCell{} & \saschaCell{} \\
KW\,31 & 27.07.--02.08. & \saschaCell{} & \saschaCell{} \\
KW\,32 & 03.08.--09.08. & \saschaCell{} & \saschaCell{} \\
KW\,33 & 10.08.--16.08. & \saschaCell{} & \saschaCell{} \\
KW\,34 & 17.08.--23.08. & \saschaCell{} & \saschaCell{} \\
KW\,35 & 24.08.--30.08. & \saschaCell{} & \saschaCell{} \\
KW\,36 & 31.08.--06.09. & \saschaCell{} & \saschaCell{} \\
KW\,37 & 07.09.--13.09. & \saschaCell{} & \saschaCell{} \\
KW\,38 & 14.09.--20.09. & \saschaCell{} & \saschaCell{} \\
KW\,39 & 21.09.--27.09. & \saschaCell{} & \saschaCell{} \\
KW\,40 & 28.09.--04.10. & \saschaCell{} & \saschaCell{} \\
\midrule
\multicolumn{3}{l}{\textbf{Summe}} & \saschaCell{} \\

\end{longtable}}

\section{Arbeitsteilung}
Das Projekt wurde von zwei Personen bearbeitet. Der vorliegende Hauptteil des Berichts beschreibt die Arbeiten von Jan Soetebeer: Scraper und Crawl-Strategie, AWS-Umgebung und Betrieb, Klassifikator, Discovery, OCR und LLM-Zweitstufe sowie die Abschlussauswertung. Der Beitrag von Sascha Lauk wird in Kapitel~\ref{ch:sascha} von ihm selbst dargestellt \saschaPH{Aufgabenbereich in zwei bis drei Sätzen}.

\section{Werkzeuge und Arbeitsweise}
Entwickelt wurde in Python~3.12 mit Git/GitHub (Branches \texttt{main}, \texttt{develop}, \texttt{feature/add-ML-Models}, Pull-Requests), Docker und der AWS-CLI. Bei der Implementierung wurde ein KI-Coding-Assistent (Claude Code) eingesetzt; die Sitzungsprotokolle dienten zusätzlich als Quelle für die Rekonstruktion des Projektverlaufs. In der LLM-Zweitstufe der Pipeline selbst kommt Claude Haiku~4.5 über Amazon Bedrock zum Einsatz (Abschnitt~\ref{sec:llm}). \TODO{Kennzeichnung des KI-Einsatzes mit den Vorgaben der FH Wedel abgleichen}

\paragraph{Umfang der Implementierung.} Das Repository umfasst rund 7.800 Zeilen Python im Crawler (\texttt{webscraper/}, 51 Dateien), rund 2.300 Zeilen im Klassifikator-Paket (\texttt{mlclassifier/}), rund 800 Zeilen in den Auswertungsskripten, rund 1.900 Zeilen Frontend-Code (HTML/CSS/JavaScript) und rund 1.800 Zeilen Infrastruktur- und Workflow-Dateien (Bootstrap-Skript, GitHub-Actions-Workflows, Terraform).
